\section{Diachronic practice: memory, anticipation, and agency inherit the same lineage}
\label{sec:diachronic_practice}

The point of the measurement-reference quotient is not confined to evidential comparison. The same
records support memory claims, anticipatory claims, and deliberative guidance. If those practices are
treated as truth-apt on the fixed scope, then they must all inherit a common referential lineage.
This section makes that dependence explicit.

\subsection{Memory and anticipation}

\begin{definition}[Memory claim]
A \emph{memory claim} is a diachronic claim whose truth conditions depend on taking a present
subject-stage to stand in the relevant same-scope lineage relation to a past record-bearing history.
\end{definition}

\begin{definition}[Anticipation claim]
An \emph{anticipation claim} is a diachronic claim whose truth conditions depend on taking a present
subject-stage to stand in the relevant same-scope lineage relation to one or more future
record-bearing histories.
\end{definition}

\begin{definition}[Referential lineage]
\label{def:referential-lineage}
For an admissible reference assignment $\rho$, the \emph{referential lineage} induced by $\rho$ is the
same-scope chain of subject-stage, successor-typing, and record-reference relations through which the
claims in $\Dset$ receive their standing-bearing values under $\rho$.
\end{definition}


\begin{definition}[Diachronic claim]
\label{def:diachronic-claim}
A \emph{diachronic claim} is any claim in the target explanatory practice whose standing-bearing value
depends on locating a present subject-stage, a past or future measurement record, or both within the
same referential lineage.
\end{definition}

\begin{theorem}[Diachronic closure]
\label{thm:diachronic-closure}
Any admissible diachronic claim factors through the same standing-preserving equivalence relation on
$\ARef$, namely the measurement-reference quotient $\ARef/\MReq$. Equivalently, there is no
admissible diachronic operator whose standing-bearing content is independent of that quotient.
\end{theorem}

\begin{proof}
By Definition~\ref{def:diachronic-claim}, a diachronic claim belongs to the target explanatory
practice only insofar as it receives a standing-bearing value from the same record-sensitive profile map
$\Prof$ used to define $\MReq$. By Theorem~\ref{thm:measurement-reference-quotient}, every
standing-bearing content space for admissible measurement-reference assignments factors through the
unique quotient $\ARef/\MReq$. So if a diachronic claim failed to factor through that quotient, two
assignments in the same $\MReq$-class would have to yield different standing-bearing values for the
claim, contradicting the definition of $\MReq$ and Theorem~\ref{thm:equivalence-collapse}. Hence
any admissible diachronic claim factors through the same standing-preserving equivalence relation.
\end{proof}

\begin{lemma}[Memory truth conditions depend on record reference]
\label{lem:memory-depends}
If two admissible assignments differ on which past record is the referent of a memory claim, then they
need not assign the same truth value to that claim.
\end{lemma}

\begin{proof}
A memory claim is true only if the present subject-stage is related, under the relevant diachronic
lineage, to the past record the claim is about. Changing which record is taken as the referent can
therefore change whether the claim is true. Hence memory truth conditions depend on record
reference.
\end{proof}

\begin{lemma}[Anticipation truth conditions depend on successor typing]
\label{lem:anticipation-depends}
If two admissible assignments differ on which future record-bearing history is the referent relevant to
an anticipation claim, then they need not assign the same truth value to that claim.
\end{lemma}

\begin{proof}
An anticipation claim attributes a future outcome to the present subject-stage. If admissible
assignments differ on which future record-bearing history is the relevant continuation, then they may
differ on whether the attributed outcome occurs in that history. So anticipation truth conditions
depend on successor typing and record reference.
\end{proof}

\subsection{Agency and deliberation}

\begin{definition}[Agential deliberation claim]
An \emph{agential deliberation claim} is a claim whose standing-bearing use guides present action by
appeal to future measurement-record outcomes treated as belonging to the same diachronic
practice-package.
\end{definition}

\begin{lemma}[Agential guidance requires typed future records]
\label{lem:agency-typed}
If a deliberative claim is to guide present action by appeal to future measurement outcomes, then the
future records appealed to must be typed as belonging to the same referential lineage as the present
subject-stage.
\end{lemma}

\begin{proof}
Deliberation links a present act to future consequences for the same diachronic practice. If the future
records appealed to are not typed as belonging to the relevant lineage, then the deliberative claim has
no determinate target for guidance. So agential guidance requires typed future records.
\end{proof}

\begin{theorem}[Decision-theoretic dependence for literal first-person anticipation]
\label{thm:decision-theoretic-dependence}
\label{prop:decision-presupposes-fixation}
Any literal first-person decision-theoretic use of branching quantum theory presupposes prior
successor typing and record reference. It does not create them.
\end{theorem}

\begin{proof}
A decision-theoretic rule can distribute weights or preferences only over already identified future
alternatives. In the present setting those alternatives are future record-bearing histories. If successor
relations and the relevant records are not already typed, then the decision rule lacks a determinate
space of future options to which its weights or preferences apply. So decision-theoretic use
presupposes rather than creates reference fixation.
\end{proof}

\begin{proposition}[Branch-aggregative retyping]
\label{prop:branch-aggregative-retyping}
A framework may avoid literal first-person anticipation by retyping deliberation into branch-aggregative
or impersonal terms. Such a framework does not remove the need for referential typing; it changes the
discourse whose truth conditions are at issue.
\end{proposition}

\begin{proof}
If a framework replaces literal first-person anticipation by impersonal aggregation over branches, then
it no longer uses the same diachronic first-person claims as the target explanatory practice. The move
therefore avoids the strongest lineage requirement only by changing the discourse at issue. That is
retyping, not elimination of the prior dependence.
\end{proof}

\subsection{A unified lineage theorem}

\begin{theorem}[Identity-preserving lineage theorem]
\label{thm:identity-preserving-lineage}
On the fixed downstream substrate, memory, anticipation, deliberation, and theory-level confirmation
all presuppose a fixed identity-preserving referential lineage through the measurement records used in
those practices.
\end{theorem}

\begin{proof}
By Lemma~\ref{lem:memory-depends}, memory claims depend on which past records are the relevant
referents. By Lemma~\ref{lem:anticipation-depends}, anticipation claims depend on successor typing
and future record reference. By Lemma~\ref{lem:agency-typed}, agential guidance depends on typed
future records. By Proposition~\ref{prop:assignment-sensitive} and
Theorem~\ref{thm:theory-level-confirmation-criterion}, theory-level confirmation depends on the
standing-bearing evidential use of the same records. These four practices therefore do not rest on four
independent structures. They rest on one identity-preserving referential lineage through the relevant
measurement records.
\end{proof}

\begin{theorem}[Diachronic inheritance theorem]
\label{thm:diachronic-inheritance}
On the fixed scope, memory, anticipation, deliberation, and theory-level confirmation all factor
through the same measurement-reference quotient. Consequently, any quotient-relevant divergence in
lineage assignment is inherited by all four practices together rather than by one of them in isolation.
\end{theorem}

\begin{proof}
By Theorem~\ref{thm:diachronic-closure}, every admissible diachronic claim factors through the unique quotient $\ARef/\MReq$; this is a special case of the general measurement-reference quotient theorem.
Memory and anticipation claims depend on which record-bearing histories are taken as their referents
(Lemmas~\ref{lem:memory-depends} and \ref{lem:anticipation-depends}); agential deliberation depends on
the same typed future records (Lemma~\ref{lem:agency-typed}); and theory-level confirmation depends on
the standing-bearing evidential use of those records (Theorem~\ref{thm:theory-level-confirmation-criterion}).
Therefore all four practices factor through the same measurement-reference quotient. If divergence is
quotient-relevant, then it is inherited by each practice through that shared factorization.
\end{proof}

\begin{theorem}[Diachronic practice dependence theorem]
\label{thm:diachronic-practice}
If admissible reference assignments remain divergent on the fixed scope in a way that changes the
referential lineage of the relevant measurement records, then memory, anticipation, deliberation, and
theory-level confirmation become assignment-relative together. No one of the four can be stabilized in
isolation while the shared quotient-level lineage remains divergent.
\end{theorem}

\begin{proof}
By Theorem~\ref{thm:identity-preserving-lineage}, the four practices share a common referential
lineage. If admissible assignments diverge in a way that changes that lineage, then by
Lemmas~\ref{lem:memory-depends} and \ref{lem:anticipation-depends} they can assign different truth
conditions to memory and anticipation claims; by Lemma~\ref{lem:agency-typed} they can change the
future records relevant to deliberation; and by Corollary~\ref{thm:confirmational-consequence} they
can change theory-level confirmation. Since the dependence is shared, the instability is shared. The
four practices therefore become assignment-relative together.
\end{proof}

\begin{corollary}[Structural priority over probabilistic treatment]
\label{cor:structural-before-probabilistic}
Questions about branch-weighting, credence, or rational choice are downstream of the lineage issue.
They presuppose, rather than settle, the fixed referential lineage required by the target explanatory
practice.
\end{corollary}

\begin{proof}
Immediate from Theorems~\ref{thm:decision-theoretic-dependence} and
\ref{thm:identity-preserving-lineage}. A probabilistic or decision-theoretic treatment can operate only
once the relevant future records and successor relations have already been typed into a fixed lineage.
\end{proof}
